% Prior Law References and Statutory Crosswalk (MAIN POINT 3). Four connecting
% movements; the dense citations live in three crosswalk tables. Every law
% resolves to its bibliography entry (eCFR, Cornell LII, or legislature URL).

This section ties H.R. 9510 to the federal and state law it builds on and to the
author's prior regulatory adaptations. Each authority resolves to its
bibliography so the reference carries the eCFR, Cornell LII, or
legislature link; the dense mapping is carried in the three crosswalk tables, and
the prose states only how the bill extends each authority without displacing it.

\subsection{Federal Law Crosswalk}

The bill rests on the existing federal framework and adds one sequencing rule to
it. Table~\ref{tab:prior-federal} sets out each authority, its citation, and the
precise way the bill builds on it. A citation correction is recorded here: the
federal antidiscrimination guarantees are codified at title 42 of the United
States Code, not title 26. Title VI appears at 42 U.S.C. § 2000d and the
Americans with Disabilities Act at 42 U.S.C. § 12101; the bill cites those titles.
An earlier informal reference to 26 U.S.C., which is the Internal Revenue Code,
for these provisions was a mislabel and is corrected throughout. Alongside the
core authorities, the bill rests on the Federal Food, Drug, and Cosmetic Act, the
Public Health Service Act, the electronic-records rule, and the investigational
pathways of 21 CFR Parts 312 and 50; the authority-and-source style of the
crosswalk follows the prior template's front matter and the national platform's
government-framework and regulatory-landscape analyses.

\begin{table}[ht]
\centering
\begin{tabularx}{\textwidth}{L{3.8cm} L{4cm} Y}
\toprule
\textbf{Authority} & \textbf{Citation} & \textbf{How the bill builds on it} \\
\midrule
Medicare \cite{usc-1395y} & 42 U.S.C. § 1395y & Medical-necessity determination stays with a licensed human \\
HIPAA \cite{cfr-hipaa} & 45 CFR Parts 160, 164 & Privacy and security of the hash-chained audit trail \\
Devices \cite{cfr-devices} & 21 CFR Parts 860 to 892 & Robot, software device doc framing \\
Title VI \cite{usc-titlevi} & 42 U.S.C. § 2000d & Bias minimization on protected characteristics \\
ADA \cite{usc-ada} & 42 U.S.C. § 12101 & Accessibility and accommodation \\
FD\&C Act \cite{usc-fdca} & 21 U.S.C. § 301 et seq. & Statutory basis for FDA authority over trial \\
PHSA \cite{usc-phsa} & 42 U.S.C. § 201 et seq. & Public-health and biologics oversight \\
Electronic records \cite{cfr-part11} & 21 CFR Part 11 & Trustworthy elec. records, signatures \\
IND pathway \cite{cfr-part312} & 21 CFR Part 312 & Trial-conduct frame the gate attaches to \\
Human subjects \cite{cfr-part50} & 21 CFR Part 50 & Human-subject protection for robots \\
\bottomrule
\end{tabularx}
\caption{Federal-law crosswalk; antidiscrimination citations corrected to 42 U.S.C.}
\label{tab:prior-federal}
\end{table}

\clearpage

\subsection{State Law Crosswalk}

Four states supply the closest analogues, each at the utilization-review or
claims layer, and the bill is designed to harmonise with and build on all four.
Table~\ref{tab:prior-state} states each requirement and the clause of this bill it
informs. New York would require submission of the algorithm and training data for
certification; Texas changed its governance between the filed version, which
required submission of the algorithm and training datasets, and the committee
substitute, which requires an annual compliance statement and the algorithms only
on a commissioner finding of noncompliance; California keeps the medical-necessity
determination with a licensed physician and bars reliance on generalized rather
than individual clinical data; and Florida bars artificial intelligence as the
sole basis for a denial and requires independent human review. The dual
state-and-federal jurisdiction treatment is from the national platform's
regulatory-landscape \cite{kawchak2026national}.

\begin{table}[ht]
\centering
\begin{tabularx}{\textwidth}{L{4.6cm} Y L{3.3cm}}
\toprule
\textbf{State and bill} & \textbf{State requirement} & \textbf{Informs which clause} \\
\midrule
New York A9149 \cite{ny-a9149} & Submit AI algorithms and training data for certification of minimized bias and evidence-based guidelines & Documentation, nondiscrimination clauses \\
Texas SB 1822 \cite{tx-sb1822} & Committee substitute: annual AI compliance statement; algorithms on a noncompliance finding & Compliance statement and enforcement trigger \\
California SB 1120 \cite{ca-sb1120} & Licensed physician retains medical necessity; no reliance on generalized data & Retained-human-review clause \\
Florida HB 527, SB 202 \cite{fl-hb527} & AI not the sole basis for denial; independent human review & Human-review backstop \\
\bottomrule
\end{tabularx}
\caption{State-law crosswalk; the Texas committee substitute is the operative version.}
\label{tab:prior-state}
\end{table}

\subsection{Prior Author Regulatory Adaptations}

The bill is consistent with, and can codify portions of, the author's prior
Physical AI regulatory adaptations. Table~\ref{tab:prior-adapt} shows how each
section-by-section adaptation maps onto a bill clause, following the
change-summary-table model of the prior template's front matter. The 21 CFR Part
312 adaptation is the structural template for this bill; the 21 CFR Part 50 and
ICH E6(R3) adaptations supply the human-subject and good-clinical-practice
content; and the national MCP-PAI standard supplies the audit, ledger, and
provenance infrastructure. The bill's own section-by-section analysis, in
Section~\ref{sec:section_by_section}, completes the mapping from adaptations
to enactable text.

\begin{table}[ht]
\centering
\begin{tabularx}{\textwidth}{L{4.25cm} Y}
\toprule
\textbf{Prior adaptation} & \textbf{How it maps onto this bill} \\
\midrule
21 CFR Part 312 \cite{kawchak2026cfr312} & Structural template; cover, style, and new-subpart model statute \\
21 CFR Part 50 \cite{kawchak2026cfr50} & Human-subject protections extended to robot-patient contact \\
ICH E6(R3) \cite{kawchak2026iche6r3, iche6r3} & Good-clinical-practice basis for human review and audit \\
National MCP-PAI \cite{kawchak2026mcp} & Hash-chained audit, ledger, and provenance infrastructure \\
\bottomrule
\end{tabularx}
\caption{Prior author adaptations mapped onto the bill.}
\label{tab:prior-adapt}
\end{table}

\subsection{Differing Governance and Corrections}

The crosswalk surfaces four distinct governance models: New York's
submit-and-certify model, the Texas attest-unless-noncompliant model, the
California physician-decision model, and the Florida human-review model. H.R.
9510 does not choose among them; it adopts a verification-before-generation
standard that sits upstream of all four, at the trial-conduct robot-code layer
they do not reach, so the states may continue to legislate at the claims layer on
top of a uniform federal conduct-layer rule. Two corrections are confirmed: the
federal antidiscrimination titles are cited at 42 U.S.C. § 2000d and 42 U.S.C.
§ 12101, and the Texas committee substitute, not the filed version, is treated as
the operative statute. The detailed mappings are in
Tables~\ref{tab:prior-federal} through~\ref{tab:prior-adapt} and are not repeated
here.
